% Part: counterfactuals
% Chapter: introduction
% Section: material-conditional

\documentclass[../../../include/open-logic-section]{subfiles}

\begin{document}

\olfileid{cnt}{int}{mat}

\olsection{বস্তুগত শর্তবচন}

ইংরেজিতে সবচেয়ে সরল রূপে শর্তবচন হলো ``If \dots then \dots''
আকারের বাক্য; এখানে \dots{}-এর প্রত্যেকটির জায়গায় আবার একটি বাক্য
থাকে। যেমন, ``পরিচারক কাজটি করে থাকলে মালী নির্দোষ।'' প্রাথমিক
যুক্তিবিদ্যার পাঠে আমরা $\lif$ সংযোজক দিয়ে এমন শর্তবচনের প্রতীকায়ন
শিখি: \dots{}-চিহ্নিত অংশগুলি, ধরা যাক, !!{formula}s $!A$ ও~$!B$
দিয়ে প্রকাশ করি; তখন পুরো শর্তবচনটি $!A \lif !B$ দিয়ে প্রকাশিত হয়।

$\lif$ সংযোজকটি \emph{সত্যমান-অপেক্ষকধর্মী}। অর্থাৎ $!A \lif !B$-এর
সত্যমান—$\True$ বা $\False$—শুধু $!A$ ও~$!B$-এর সত্যমান দিয়ে
নির্ধারিত হয়: $!A$ মিথ্যা অথবা $!B$ সত্য হলেই $!A \lif !B$ সত্য;
অন্যথায় মিথ্যা। সত্যমান-আরোপ~$\pAssign v$-এর সাপেক্ষে আমরা
$\pSat{v}{!A \lif !B}$ সংজ্ঞায়িত করি তখন এবং কেবল তখনই, যখন
$\pSat/{v}{!A}$ অথবা $\pSat{v}{!B}$। এই অর্থতত্ত্বে $\lif$
সংযোজককে \emph{বস্তুগত শর্তবচন} বলা হয়।

এই সংজ্ঞা থেকে কয়েকটি প্রাথমিক যৌক্তিক তথ্য পাওয়া যায়। প্রথমত,
বস্তুগত শর্তবচনের জন্য নিঃসরণ উপপাদ্যটি সত্য:
\begin{align}
  \text{যদি } \Gamma, !A \Entails !B \text{ হয়, তবে } \Gamma \Entails !A \lif !B
\end{align}
এটি সত্যমান-অপেক্ষকধর্মী: $!A \lif !B$ ও
$\lnot !A \lor !B$ সমতুল্য:
\begin{align}
  !A \lif !B & \Entails \lnot !A \lor !B\\
  \lnot !A \lor !B & \Entails !A \lif !B
  \intertext{অনুপক্ষ থেকে এবং পূর্বপক্ষের নঞর্থকরণ থেকে
    বস্তুগত শর্তবচনটি অনুসৃত হয়:}
  !B & \Entails !A \lif !B\\
  \lnot !A & \Entails !A \lif !B
  \intertext{মিথ্যা বস্তুগত শর্তবচন তার পূর্বপক্ষ ও অনুপক্ষের
    নঞর্থকরণের সংযোজনের সমতুল্য: $!A \lif !B$ মিথ্যা হলে
    $!A \land \lnot !B$ সত্য, এবং বিপরীতটিও সত্য:}
  \lnot(!A \lif !B) & \Entails !A \land \lnot !B\\
  !A \land \lnot !B & \Entails \lnot(!A \lif !B)
  \intertext{বস্তুগত শর্তবচনের জন্য মোডাস পোনেন্স বিধি চলে:}
  !A, !A \lif !B & \Entails !B
  \intertext{বস্তুগত শর্তবচনে দুটি অনুপক্ষ একত্র করা যায়:}
  !A \lif !B,  !A \lif !C & \Entails !A \lif (!B \land !C)
  \intertext{পূর্বপক্ষকে সব সময় শক্তিশালী করা যায়; অর্থাৎ
    শর্তবচনটির \emph{সত্যতা পূর্বপক্ষ শক্তিশালী করলেও সংরক্ষিত}:}
  !A \lif !B & \Entails (!A \land !C) \lif !B
  \intertext{বস্তুগত শর্তবচন পরিযায়ী; অর্থাৎ শৃঙ্খল-বিধিটি বৈধ:}
  !A \lif !B, !B \lif !C & \Entails !A \lif !C
  \intertext{বস্তুগত শর্তবচন তার প্রতিবিপরীত রূপের সমতুল্য:}
  !A \lif !B & \Entails \lnot !B \lif \lnot !A\\
  \lnot !B \lif \lnot !A & \Entails !A \lif !B
\end{align}

গাণিতিক যুক্তিবিচারে এই সব অনুসিদ্ধান্তই কার্যকর এবং নির্দ্বিধায়
গ্রহণযোগ্য। কিন্তু দর্শন ও ভাষাবিজ্ঞানের সাহিত্যে অ-গাণিতিক প্রসঙ্গে
এগুলির অনুরূপ অনুসিদ্ধান্তের বিরুদ্ধে আপাত-প্রতিদৃষ্টান্ত প্রচুর।
সেগুলি থেকে বোঝা যায়, ইংরেজি ``if \dots then \dots'' বক্তব্য
প্রতীকায়নের জন্য বস্তুগত শর্তবচন~$\lif$ যথাযথ সংযোজক নয়—অন্তত
সব ক্ষেত্রে যথাযথ নয়।

\end{document}
